\chapter{Achilles Tendonitis}

\section*{Definition and Overview}
Achilles tendonitis refers to pain, swelling, and irritation of the Achilles tendon, the thick band of tissue that connects the calf muscles to the heel bone. Health professionals increasingly prefer the term ``Achilles tendinopathy'', because the underlying problem is usually degeneration and failed healing rather than active inflammation; however, both names describe the same common overuse condition.

There are two main forms. Non-insertional (mid-portion) tendinopathy affects the middle part of the tendon, a few centimetres above the heel, and is the more common type. Insertional tendinopathy occurs where the tendon attaches to the heel bone, is often associated with a bony prominence called a Haglund's deformity, and tends to be more stubborn to treat. This chapter covers both forms.

\section*{Epidemiology}
Achilles tendinopathy is one of the most common running injuries and is seen frequently in sports that involve running, jumping, and rapid changes of direction, including athletics, football, netball, and basketball. It affects men and women alike, with the highest rates among active adults in their 30s to 50s.

It also occurs in people who are not athletes, particularly those who are overweight or who suddenly increase their activity level. The risk rises with age as the tendon becomes less elastic and slower to repair, and it is more common in people with conditions such as diabetes, gout, and inflammatory arthritis.

\section*{Aetiology and Pathophysiology}
Achilles tendinopathy is usually caused by overuse: the tendon is loaded repeatedly beyond its capacity to repair, so microscopic collagen damage accumulates faster than it can heal. Contributing factors include a sudden increase in training distance or intensity, inadequate warm-up, tight or weak calf muscles, poor or worn footwear, and running on hard or uneven surfaces. Biomechanical factors such as flat feet, high arches, or stiff ankles can increase strain on the tendon, and certain medicines, including fluoroquinolone antibiotics and corticosteroids, are associated with tendon injury.

At the microscopic level, the tendon becomes thickened and disorganised, with abnormal collagen fibres, increased fluid, and small areas of degeneration. Rather than being a purely inflammatory condition, it is best understood as a failed healing response; this is why it is often slow to resolve and prone to recurrence.

\section*{Clinical Features}
\begin{itemize}
    \item \textbf{Pain and aching:} at the back of the heel or along the tendon, worse with activity
    \item \textbf{Morning stiffness:} the tendon feels tight and stiff on first standing, easing as it warms up
    \item \textbf{Start-up pain:} discomfort at the beginning of exercise that may ease, then return afterwards
    \item \textbf{Swelling and thickening:} the tendon may feel swollen, warm, and tender to press
    \item \textbf{Crepitus:} a creaking or grating sensation when moving the ankle
    \item \textbf{Functional limitation:} pain when walking up hills or stairs, or after sitting for a long time
\end{itemize}

\section*{Differential Diagnosis}
Pain around the back of the ankle and heel has several possible causes that must be distinguished.

\begin{itemize}
    \item Acute rupture of the Achilles tendon
    \item Plantar fasciitis (pain along the underside of the heel)
    \item Retrocalcaneal bursitis (inflammation of the bursa behind the heel)
    \item Calf muscle strain or tear
    \item Posterior tibial tendon dysfunction (inner side of the ankle)
    \item Ankle joint arthritis or a stress fracture of the heel bone
    \item Haglund's deformity (bony prominence at the back of the heel)
\end{itemize}

\section*{When to Seek Medical Care}
See a doctor or physiotherapist if heel or tendon pain lasts for more than a few weeks despite rest, if it is severe, or if it is stopping you from walking normally. Seek urgent assessment if you hear or feel a sudden snap at the back of the ankle with sharp pain and difficulty pushing off, as this suggests a rupture rather than tendonitis.

\textbf{Call 000 (or local emergency services) for:}
\begin{itemize}
    \item Sudden severe pain after a major injury, with an inability to stand or bear weight
    \item The foot or leg below the injury turning cold, pale, or numb
    \item An open wound or visible bone after trauma to the ankle
    \item Any sudden, severe, or concerning symptom
\end{itemize}

\section*{Diagnosis and Investigations}
\begin{itemize}
    \item \textbf{History:} questions about activity levels, training changes, footwear, and how the pain started
    \item \textbf{Examination:} assessment of tenderness, swelling, and thickening along the tendon, and of ankle movement and strength
    \item \textbf{Thompson (calf squeeze) test:} if squeezing the calf does not move the foot, the tendon is likely torn
    \item \textbf{Ultrasound:} shows thickening, tears, and blood flow within the tendon, and can guide injections
    \item \textbf{MRI:} provides detailed images when the diagnosis is uncertain or surgery is being considered
    \item \textbf{X-ray:} used mainly to exclude bony problems such as fractures, spurs, or a Haglund's deformity at the heel
\end{itemize}

\section*{Management}
\begin{itemize}
    \item \textbf{Activity modification:} reduce or change activities that aggravate the tendon, such as hill running or sprinting
    \item \textbf{Ice:} short applications in the early stages to ease pain and swelling
    \item \textbf{Pain relief:} short-term paracetamol or non-steroidal anti-inflammatory drugs, used on advice
    \item \textbf{Physiotherapy:} the mainstay of treatment, especially graduated eccentric (lengthening) strengthening exercises
    \item \textbf{Footwear and orthoses:} supportive shoes, heel lifts, or orthotic insoles can reduce strain
    \item \textbf{Shockwave therapy:} may help cases that have not improved with exercise
    \item \textbf{Injections:} considered cautiously, as some agents can weaken the tendon
    \item \textbf{Surgery:} rarely needed for persistent cases, and always followed by rehabilitation
\end{itemize}

\section*{Complications}
\begin{itemize}
    \item Progression to a partial or complete rupture if the tendon is loaded heavily while injured
    \item Chronic, long-lasting pain that limits sport and daily life
    \item Permanent thickening of the tendon
    \item Weakening of the tendon from repeated corticosteroid injections
    \item Complications of surgery, including infection and nerve injury
    \item Overuse of the other leg as people shift weight to avoid discomfort
\end{itemize}

\section*{Prognosis and Outcomes}
Most people improve with physiotherapy, activity modification, and patience, but recovery is often slow and commonly takes several months. Symptoms may ease and then return, particularly if training is increased too quickly, and the tendon may remain thickened even after pain resolves. The outlook is generally better for mid-portion tendinopathy than for insertional problems, which can be harder to treat. A structured rehabilitation programme offers the best chance of returning to sport without recurrence, and full recovery is achievable in the majority of cases.

\section*{Prevention and Self-Care}
\begin{itemize}
    \item Build up training gradually, avoiding sudden increases in distance, intensity, or hill running
    \item Warm up properly and include regular calf stretches in your routine
    \item Strengthen the calf muscles regularly with slow, controlled exercises
    \item Wear supportive, well-fitted footwear and replace worn shoes
    \item Cross-train with low-impact activities to reduce repetitive load
    \item Treat early symptoms promptly rather than ``running through'' the pain
\end{itemize}

\section*{Sources and Further Reading}
The following resources provide reliable, up-to-date information on Achilles tendonitis and tendinopathy.

\begin{itemize}
    \item NHS. \textit{Information on Achilles tendon problems, including causes and self-care.} https://www.nhs.uk/conditions/
    \item Mayo Clinic. \textit{Guide to Achilles tendinitis, diagnosis, and treatment.} https://www.mayoclinic.org/
    \item MSD Manual. \textit{Professional and patient education on tendon disorders.} https://www.msdmanuals.com/
    \item MedlinePlus. \textit{Health information on Achilles tendon injuries.} https://medlineplus.gov/
    \item NICE. \textit{Clinical knowledge and guidance on tendon disorders and rehabilitation.} https://www.nice.org.uk/
\end{itemize}
